
\clearpage
\chapter{ארכיטקטורות \textenglish{Transformer} לטורי זמן}

\section{\textenglish{PatchTST}}

\textenglish{PatchTST} \cite{nie2022patchtst} מגדיר מחדש את ייצוג הקלט: במקום להזין ערך בודד לכל צעד זמן, המודל מחלק את הרצף לחלונות חופפים (\textenglish{patches}) ומייצג כל חלון כטוקן יחיד. גישה זו מפחיתה את אורך הרצף לאחר הפלסה ב-\textenglish{patch size} ומאפשרת ל-\textenglish{attention} לתפוס תבניות גלובליות בעלות ייצוג מקומי פלוסי.

\section{\textenglish{Autoformer}}

\textenglish{Autoformer} \cite{wu2021autoformer} מציג מנגנון \textenglish{Auto-Correlation} המחליף את ה-\textenglish{self-attention} הסטנדרטי:

\begin{equation}
\mathcal{R}_{QK}(\tau) = \lim_{L\to\infty} \frac{1}{L} \sum_{t=1}^L Q(t) \cdot K(t-\tau)
\end{equation}

מנגנון זה מנצל עצמי-מתאם (\textenglish{autocorrelation}) בין הקוורי לבין המפתח ומסיר מגבלת הזיכרון הריבועית, תוך שמירה על תפיסת תלויות זמניות.

\section{\textenglish{FEDformer}}

\textenglish{FEDformer} \cite{zhou2022fedformer} מוסיף פירוק תדר (\textenglish{frequency domain decomposition}) לצינור ה-\textenglish{Transformer}. על ידי עבודה בתחום פורייה ובחירת מרכיבי תדר מובהקים בלבד, המודל מפחית את מורכבות ה-\textenglish{attention} מ-\textenglish{O(n²)} ל-\textenglish{O(n log n)}.
